% R7 regime diagram: the deterrence frontier rho = G/(dB) at G=10, d=0.8,
% in the (bond B, audit rate rho) plane, B in [10,100]. Curve values
% recomputed directly from rho_curve = clip(G/(d*B_range), 0, 1) in
% whitepaper/research/figures/src/r7_figures.py, lines 97-102 (sampled
% at 13 representative B values spanning the plotted range, including
% the B=12.5 point where the curve first drops below the rho=1 clip
% ceiling). Marked point (50, 0.25) and the tower inset text are
% transcribed from the script's ax.plot / ax.annotate / ax.text calls.
% [verified: script r7_figures.py]
%
% REVIEW FIXES (exposition review B2), neither of which touches the data:
%   (1) the "cheating pays" region was a seagreen!12 / shipred!8 tint pair
%       -- two very light washes that are near-indistinguishable in
%       greyscale, leaving the legend to carry the whole load. It is now a
%       hatch, so the two half-planes differ in texture as well as hue.
%       (\usetikzlibrary{patterns} is loaded by paper3.tex.)
%   (2) the inset and the caption attributed the decay rate (1-rho*d)=0.8
%       to C=8. The rate is independent of C -- it is identical at C=1.
%       What C sets is where the geometric decay STARTS (the threshold
%       G_k <= C*B). Both now say so.
\begin{figure}[htbp]\centering
\begin{tikzpicture}
\begin{axis}[
  harbor curve,
  width=0.78\textwidth,height=6.2cm,
  xlabel={bond size $B$},
  ylabel={audit rate $\rho$},
  title={R7 regime --- the deterrence frontier $\rho d B = G$},
  xmin=10,xmax=100,ymin=0,ymax=1,
  legend pos=north east,
]
\addplot[fill=seagreen!12,draw=none] coordinates {
  (10,1.0) (12.5,1.0) (15,0.833333) (20,0.625) (25,0.5) (30,0.416667)
  (40,0.3125) (50,0.25) (60,0.208333) (70,0.178571) (80,0.15625)
  (90,0.138889) (100,0.125) (100,1.0)
} \closedcycle;
\addlegendentry{deterrence holds (cheating unprofitable)}
\addplot[pattern=north east lines,pattern color=shipred!45,draw=none] coordinates {
  (10,1.0) (12.5,1.0) (15,0.833333) (20,0.625) (25,0.5) (30,0.416667)
  (40,0.3125) (50,0.25) (60,0.208333) (70,0.178571) (80,0.15625)
  (90,0.138889) (100,0.125) (100,0)
} \closedcycle;
\addlegendentry{cheating pays}
\addplot[harborblue,thick] coordinates {
  (10,1.0) (12.5,1.0) (15,0.833333) (20,0.625) (25,0.5) (30,0.416667)
  (40,0.3125) (50,0.25) (60,0.208333) (70,0.178571) (80,0.15625)
  (90,0.138889) (100,0.125)
};
\addlegendentry{$\rho = G/(dB)$}
\addplot[only marks,mark=*,mark size=2.2pt,shipred,forget plot] coordinates {(50,0.25)};
\node[regimebox,anchor=south west,align=left,font=\scriptsize,draw=shipred,fill=white]
  (measnote) at (axis cs:65,0.05)
  {measured: cheat payoff\\ exactly $0$ at $\rho^\star$\\ {[internal, \texttt{b2\_tower.py}]}};
\draw[-{Stealth[length=2mm]},shipred,thin] (axis cs:50,0.25) -- (measnote.north);
\node[regimebox,anchor=north west,align=left,font=\scriptsize,draw=harborblue,fill=yellow!15]
  at (rel axis cs:0.03,0.60)
  {tower: once $G_k \le CB$, corrupt value\\ decays $(1-\rho d)^k = 0.8^k$ per level\\ (rate independent of $C$;\\ $C$ sets where the decay starts)};
\end{axis}
\end{tikzpicture}
\caption{The regime diagram: the deterrence frontier $\rho d B = G$ in the (bond $B$, audit rate $\rho$) plane at $G{=}10$, $d{=}0.8$. Above the curve $\rho = G/(dB)$, cheating is unprofitable; below it, cheating pays. The marked point is the running example $(B,\rho^\star) = (50, 0.25)$, where the measured cheat payoff is exactly zero [internal, \texttt{b2\_tower.py}]. Below the curve, the region is hatched rather than tinted so the two half-planes stay distinguishable in greyscale. The inset records the tower consequence of \S\ref{sec:tower}: below the bribery threshold $G_k \le CB$, surviving corrupt value decays as $(1-\rho d)^k$ --- a $\times 0.8$-per-level collapse at these parameters, at every $C$. What $C$ controls is not the rate but where the decay begins.}
\label{fig:regime}
\end{figure}
